\phantomsection
\addcontentsline{toc}{section}{References}
\begin{thebibliography}{99}
\raggedright
\bibitem{Miyao}
T. Miyao, \emph{Upper bounds on the charge susceptibility of many-electron
systems coupled to the quantized radiation field},
Letters in Mathematical Physics \textbf{105} (2015), 1119--1133.
\href{https://doi.org/10.1007/s11005-015-0775-9}{doi:10.1007/s11005-015-0775-9};
\href{https://arxiv.org/abs/1405.0112}{arXiv:1405.0112}.

\bibitem{LS}
E. H. Lieb and H. Siedentop,
\emph{Renormalization of the Regularized Relativistic Electron-Positron Field},
Communications in Mathematical Physics \textbf{213} (2000), 673--683.
\href{https://doi.org/10.1007/s002200000265}{doi:10.1007/s002200000265};
\href{https://arxiv.org/abs/math-ph/0003001}{arXiv:math-ph/0003001}.

\bibitem{HLS}
C. Hainzl, M. Lewin, and \'{E}. S\'er\'e,
\emph{Existence of a Stable Polarized Vacuum in the
Bogoliubov--Dirac--Fock Approximation},
Communications in Mathematical Physics \textbf{257} (2005), 515--562.
\href{https://doi.org/10.1007/s00220-005-1343-4}{doi:10.1007/s00220-005-1343-4};
\href{https://arxiv.org/abs/math-ph/0403005}{arXiv:math-ph/0403005}.

\bibitem{GLS}
P. Gravejat, M. Lewin, and \'{E}. S\'er\'e,
\emph{Renormalization and asymptotic expansion of Dirac's polarized vacuum},
Communications in Mathematical Physics \textbf{306} (2011), 1--33.
\href{https://doi.org/10.1007/s00220-011-1271-4}{doi:10.1007/s00220-011-1271-4};
\href{https://arxiv.org/abs/1004.1734}{arXiv:1004.1734}.

\bibitem{GHLS}
P. Gravejat, C. Hainzl, M. Lewin, and \'{E}. S\'er\'e,
\emph{Construction of the Pauli--Villars-Regulated Dirac Vacuum
in Electromagnetic Fields},
Archive for Rational Mechanics and Analysis \textbf{208} (2013), 603--665.
\href{https://doi.org/10.1007/s00205-012-0609-1}{doi:10.1007/s00205-012-0609-1};
\href{https://arxiv.org/abs/1204.2893}{arXiv:1204.2893}.

\bibitem{BBHS}
V. Bach, J.-M. Barbaroux, B. Helffer, and H. Siedentop,
\emph{Stability of Matter for the Hartree--Fock Functional of the
Relativistic Electron-Positron Field},
Documenta Mathematica \textbf{3} (1998), 353--364.
\href{https://doi.org/10.4171/DM/51}{doi:10.4171/DM/51}.
\bibitem{FSS}
J. Fr\"ohlich, B. Simon, and T. Spencer,
\emph{Infrared bounds, phase transitions and continuous symmetry breaking},
Communications in Mathematical Physics \textbf{50} (1976), 79--95.
\href{https://doi.org/10.1007/BF01608557}{doi:10.1007/BF01608557}.

\bibitem{Seiler}
E. Seiler, \emph{Upper bound on the color-confining potential},
Physical Review D \textbf{18} (1978), 482--483.
\href{https://doi.org/10.1103/PhysRevD.18.482}{doi:10.1103/PhysRevD.18.482}.

\bibitem{Wei}
Z. C. Wei, C. Wu, Y. Li, S. Zhang, and T. Xiang,
\emph{Majorana Positivity and the Fermion Sign Problem of Quantum Monte Carlo Simulations},
Physical Review Letters \textbf{116} (2016), 250601.
\href{https://doi.org/10.1103/PhysRevLett.116.250601}{doi:10.1103/PhysRevLett.116.250601};
\href{https://arxiv.org/abs/1601.01994}{arXiv:1601.01994}.

\bibitem{LiJiangYao}
Z.-X. Li, Y.-F. Jiang, and H. Yao,
\emph{Majorana-Time-Reversal Symmetries: A Fundamental Principle for
Sign-Problem-Free Quantum Monte Carlo Simulations},
Physical Review Letters \textbf{117} (2016), 267002.
\href{https://doi.org/10.1103/PhysRevLett.117.267002}{doi:10.1103/PhysRevLett.117.267002};
\href{https://arxiv.org/abs/1601.05780}{arXiv:1601.05780}.

\bibitem{VafaWitten}
C. Vafa and E. Witten,
\emph{Restrictions on symmetry breaking in vector-like gauge theories},
Nuclear Physics B \textbf{234} (1984), 173--188.
\href{https://doi.org/10.1016/0550-3213(84)90230-X}{doi:10.1016/0550-3213(84)90230-X}.

\bibitem{KogutSusskind}
J. Kogut and L. Susskind,
\emph{Hamiltonian formulation of Wilson's lattice gauge theories},
Physical Review D \textbf{11} (1975), 395--408.
\href{https://doi.org/10.1103/PhysRevD.11.395}{doi:10.1103/PhysRevD.11.395}.

\bibitem{Zache}
T. V. Zache, F. Hebenstreit, F. Jendrzejewski, M. K. Oberthaler,
J. Berges, and P. Hauke,
\emph{Quantum simulation of lattice gauge theories using Wilson fermions},
Quantum Science and Technology \textbf{3} (2018), 034010.
\href{https://doi.org/10.1088/2058-9565/aac33b}{doi:10.1088/2058-9565/aac33b};
\href{https://arxiv.org/abs/1802.06704}{arXiv:1802.06704}.
\bibitem{EKM}
A. X. El-Khadra, A. S. Kronfeld, and P. B. Mackenzie,
\emph{Massive Fermions in Lattice Gauge Theory},
Physical Review D \textbf{55} (1997), 3933--3957.
\href{https://doi.org/10.1103/PhysRevD.55.3933}{doi:10.1103/PhysRevD.55.3933};
\href{https://arxiv.org/abs/hep-lat/9604004}{arXiv:hep-lat/9604004}.

\bibitem{BDL}
S. Bravyi, D. P. DiVincenzo, and D. Loss,
\emph{Schrieffer--Wolff transformation for quantum many-body systems},
Annals of Physics \textbf{326} (2011), 2793--2826.
\href{https://doi.org/10.1016/j.aop.2011.06.004}{doi:10.1016/j.aop.2011.06.004};
\href{https://arxiv.org/abs/1105.0675}{arXiv:1105.0675}.

\bibitem{Fry}
M. P. Fry, \emph{Stability of QED},
Physical Review D \textbf{84} (2011), 065021.
\href{https://doi.org/10.1103/PhysRevD.84.065021}{doi:10.1103/PhysRevD.84.065021};
\href{https://arxiv.org/abs/1108.4345v2}{arXiv:1108.4345v2}.
The two assumptions identified in Section IV of the revised manuscript
are retained in the discussion here.

\bibitem{HLSThermoI}
C. Hainzl, M. Lewin, and J. P. Solovej,
\emph{The thermodynamic limit of quantum Coulomb systems. Part I.
General theory},
Advances in Mathematics \textbf{221} (2009), 454--487.
\href{https://doi.org/10.1016/j.aim.2008.12.010}{doi:10.1016/j.aim.2008.12.010};
\href{https://arxiv.org/abs/0806.1708}{arXiv:0806.1708}.

\bibitem{HLSThermoII}
C. Hainzl, M. Lewin, and J. P. Solovej,
\emph{The thermodynamic limit of quantum Coulomb systems. Part II.
Applications},
Advances in Mathematics \textbf{221} (2009), 488--546.
\href{https://doi.org/10.1016/j.aim.2008.12.011}{doi:10.1016/j.aim.2008.12.011};
\href{https://arxiv.org/abs/0806.1709}{arXiv:0806.1709}.

\bibitem{LY1988}
E. H. Lieb and H.-T. Yau,
\emph{The Stability and Instability of Relativistic Matter},
Communications in Mathematical Physics \textbf{118} (1988), 177--213.
\href{https://doi.org/10.1007/BF01218577}{doi:10.1007/BF01218577}.

\bibitem{LL2002}
E. H. Lieb and M. Loss,
\emph{Stability of a Model of Relativistic Quantum Electrodynamics},
Communications in Mathematical Physics \textbf{228} (2002), 561--588.
\href{https://doi.org/10.1007/s002200200665}{doi:10.1007/s002200200665};
\href{https://arxiv.org/abs/math-ph/0109002}{arXiv:math-ph/0109002}.

\bibitem{Lieb2002}
E. H. Lieb, \emph{The Stability of Matter and Quantum Electrodynamics},
Heisenberg Symposium contribution (2002).
\href{https://arxiv.org/abs/math-ph/0209034}{arXiv:math-ph/0209034}.

\bibitem{JW2018}
C. D. J\"akel and W. F. Wreszinski,
\emph{Stability of relativistic quantum electrodynamics in the Coulomb gauge},
Journal of Mathematical Physics \textbf{59} (2018), 032303.
\href{https://doi.org/10.1063/1.5011031}{doi:10.1063/1.5011031};
\href{https://arxiv.org/abs/1705.04652}{arXiv:1705.04652}.

\bibitem{FLS2007}
R. L. Frank, E. H. Lieb, and R. Seiringer,
\emph{Stability of relativistic matter with magnetic fields for nuclear
charges up to the critical value},
Communications in Mathematical Physics \textbf{275} (2007), 479--489.
\href{https://doi.org/10.1007/s00220-007-0307-2}{doi:10.1007/s00220-007-0307-2};
\href{https://arxiv.org/abs/math-ph/0610062}{arXiv:math-ph/0610062}.
\end{thebibliography}
